\section{Introduction: Tracing Search to Evidence}
\label{sec:introduction}
\IEEEPARstart{J}{arambit} studies computational design for a titanium hooked utility knife with an integral ring. Its largest admitted cohort contains $1{,}954{,}560$ evaluated birth records across ten completed runs. Smaller historical and screening collections supply different observations, and separately assessed exports connect particular geometries to numerical records. A candidate's search context and its assessment history therefore need to be read together.

We use \emph{Engineered Intelligence} for a process connecting a versioned representation and decoder, recorded variation and selection, declared objectives, and inspectable observations. Jarambit's documented schema has 128 numeric coordinates, including conditional and discretized parameters, plus a categorical material field. A genome, a contextual search score, and an assessed artifact each describe a different part of that process.

\subsection{Research Question and Three Findings}
\label{subsec:research-question-and-results}
The study asks: \emph{what can the retained links establish about a computational design result?} Three findings organize the answer:
\begin{enumerate}
  \item \textbf{Coverage depends on the observation stage.} Extensive genome and score tables coexist with selective parentage, geometry and later-assessment records. The historical genealogy, native populations and intermediate screens retain different kinds of evidence.
  \item \textbf{Search rank and a recorded assessment predicate describe different outcomes.} In historical Run~1, the highest-scoring birth, Individual~1095, fails that run's recorded certification predicate, while lower-ranked Individual~1543 meets it. This categorical comparison is scoped to the earlier predicate; later all-requirement qualification is a separate assessment.
  \item \textbf{Numerical agreement and error estimation answer different questions.} Selected solid-model observables agree closely between two implementations on a shared mesh. Separate plate-refinement sequences refuse the chosen extrapolated-error estimates. Each result carries its own inputs and inference scope.
\end{enumerate}

\subsection{Contribution and Reading Guide}
\label{subsec:contributions}
The contribution is a source-linked case study joining representation semantics, candidate identity, score context and candidate-specific observations in one inspectable account. Recorded ancestry illustrates reuse within the fixed representation. Coverage views distinguish available observations from missing ones, and the numerical example separates shared-input consistency from mesh-study admission. These concrete links let a reader trace what a displayed result means without treating a collection of scores as a uniformly assessed population.

Section~\ref{sec:study-design} defines the records and counting units; Sections~\ref{sec:genome}--\ref{sec:evolution} explain representation and inheritance. Section~\ref{sec:results} develops the findings, and Section~\ref{sec:verification} interprets the numerical evidence. The analysis is computational and retrospective; Section~\ref{sec:discussion} consolidates its scope and limits.
